\chapter{Optimization Framework} 
\label{chap:optimization-framework}

% Chapter introduction: written last.

\section{Risk-Based Portfolio Optimization}\label{sec:risk-based-optimization}

Portfolio optimization theory studies how an investor should allocate their wealth across a set of different risky assets. In one of the most 
influential works in finance, Markowitz \cite{Markowitz1952} proposes to choose among the assets by trading off the expected
portfolio return against its variance. This requires estimates of the expected asset returns, which are
less precise than estimates of variances and covariances. Merton \cite{Merton1980} shows that the precision
of an estimated expected return depends only on the length of the sample period and not on how often the returns are
observed within it, whereas the precision of an estimated variance increases with the observation frequency. Michaud
\cite{Michaud1989} describes mean-variance optimizers as estimation-error maximizers, since they overweight the assets
with large estimated expected returns and small estimated variances, which are the assets most likely to have large
estimation errors. DeMiguel et al.~\cite{DeMiguelGarlappiUppal2009} find that, out of sample, none of the optimized
portfolios they consider consistently outperforms the equally weighted portfolio, because the gain from optimal
diversification is more than offset by estimation error. Following Lucescu and Ulrych
\cite{RegularizedMultiCurrencyExpectedShortfallPortfolios}, we therefore first consider portfolios that minimize a
measure of risk alone and require no estimates of expected returns. We consider the expected portfolio return later in
this chapter.

We initially formulate the portfolio optimization problem for an investor who holds $N$ assets denominated in a single currency, and extend the setting
to multiple currencies in the next section. Let $\bm{x}_t \in \R^N$ denote the portfolio weights chosen at time $t$ and
held until $t+1$, and $\bm{R}_{t+1} \in \R^N$ the vector of simple asset returns over the period from $t$ to $t+1$. The
portfolio return at time $t+1$ is then $\bm{x}_t^\top \bm{R}_{t+1}$, and the weights satisfy the budget constraint
$\bm{x}_t^\top \bm{1}_N = 1$, where $\bm{1}_N \in \R^N$ is the vector of ones. A risk measure $\rho$ assigns to the
random portfolio return a real number that quantifies its risk. We write
$\rho$ as a function of the portfolio return $\bm{x}_t^\top \bm{R}_{t+1}$, so that the risk-based portfolio optimization problem is
\begin{equation} \label{eq:risk-based-problem}
    \bm{x}_t^* = \argmin_{\bm{x}_t} \rho\bigl( \bm{x}_t^\top \bm{R}_{t+1} \bigr), \quad \text{s.t. } \bm{x}_t^\top \bm{1}_N = 1.
\end{equation}
The risk measure is evaluated under the distribution of $\bm{R}_{t+1}$ conditional on the information available at time
$t$.

\subsection{Variance}\label{sec:risk-variance}

In the portfolio theory of Markowitz \cite{Markowitz1952}, the risk of a portfolio is measured by the variance of its
return. Let $\mathbb{E}_t[\cdot]$ denote the expectation conditional on the information available at time $t$, and let
$\bm{\mu}_t = \mathbb{E}_t[\bm{R}_{t+1}] \in \R^N$ and
$\bm{\Sigma}_t = \mathbb{E}_t\bigl[ (\bm{R}_{t+1} - \bm{\mu}_t)(\bm{R}_{t+1} - \bm{\mu}_t)^\top \bigr] \in \R^{N \times N}$
denote the expected asset return and the conditional covariance matrix of the asset returns. The variance of the portfolio
return is
\begin{equation*}
    \text{Var}_t\bigl( \bm{x}_t^\top \bm{R}_{t+1} \bigr)
    = \mathbb{E}_t\bigl[ (\bm{x}_t^\top \bm{\mu}_t - \bm{x}_t^\top \bm{R}_{t+1})^2 \bigr]
    = \bm{x}_t^\top \bm{\Sigma}_t \bm{x}_t.
\end{equation*}
The variance thus depends on the distribution of the asset returns only through their covariance matrix.

With the variance as the risk measure, the objective of problem \eqref{eq:risk-based-problem} is the quadratic form
$\bm{x}_t^\top \bm{\Sigma}_t \bm{x}_t$. Since $\bm{\Sigma}_t$ is positive semi-definite, the objective is convex in the
weights, and the problem is a convex quadratic program, which standard solvers handle efficiently
\cite{boyd2004convex}. If $\bm{\Sigma}_t$ is positive definite, the first-order conditions of its Lagrangian further give
the unique solution in closed form
\begin{equation*}
    \bm{x}_t^* = \frac{\bm{\Sigma}_t^{-1} \bm{1}_N}{\bm{1}_N^\top \bm{\Sigma}_t^{-1} \bm{1}_N},
\end{equation*}
which is known as the global minimum-variance portfolio. The closed form is a particularity of the variance, since the
problems under the other risk measures we consider have no closed-form solution in general and must be solved
numerically. It no longer applies once we add the leverage constraint and other penalties later in this chapter, although
the problem remains a convex quadratic program. The closed form also shows that the weights depend on the covariance
matrix through its inverse, so that errors in an estimate of $\bm{\Sigma}_t$ carry over directly to the portfolio.

Since the conditional distribution of $\bm{R}_{t+1}$ is unknown, $\bm{\Sigma}_t$ has to be estimated. For the variance
and for the other risk measures later in the chapter, we use $q$ historical realizations $\bm{y}_1, \dots, \bm{y}_q \in \R^N$ of
$\bm{R}_{t+1}$. In the empirical analysis, these are the daily asset returns over a lookback window of the ten years
before time $t$, so that $q \approx 2{,}500$. We denote their sample mean
by $\bar{\bm{y}} = \frac{1}{q} \sum_{k=1}^q \bm{y}_k$.

The natural estimate of $\bm{\Sigma}_t$ is the sample covariance matrix of the realizations
\begin{equation*}
    \bar{\bm{\Sigma}}_t = \frac{1}{q} \sum_{k=1}^q (\bm{y}_k - \bar{\bm{y}})(\bm{y}_k - \bar{\bm{y}})^\top.
\end{equation*}
It has $N(N+1)/2$ distinct
entries, all estimated from the same $q$ observations, and when $N$ is large relative to $q$, its eigenvalues are more
dispersed than those of $\bm{\Sigma}_t$, with the smallest eigenvalues biased downward \cite{LedoitWolf2004}. Through
the inverse in the closed form, the minimum-variance portfolio then concentrates on the combinations of assets whose
risk is most underestimated, and its in-sample variance understates its out-of-sample variance \cite{ElKaroui2010}.
When $N \ge q$, the sample covariance matrix is singular, and the closed form does not exist.

We therefore estimate $\bm{\Sigma}_t$ with the shrinkage estimator of Ledoit and Wolf \cite{LedoitWolf2003}, which
combines the sample covariance matrix with a structured target. The sample covariance matrix imposes no structure but
has a large estimation error when $N$ is large relative to $q$, while a structured target has few parameters and
therefore a small estimation error but is biased whenever its structure does not hold. A weighted average of the two
trades off the estimation error of the one against the bias of the other. Ledoit and Wolf use the constant-correlation
target, which keeps the sample variances of the assets and replaces every pairwise correlation by the average sample
correlation.

To define the estimator, let $y_{k,i}$ denote the $i$th entry of $\bm{y}_k$, $\bar{y}_i$ the $i$th entry of
$\bar{\bm{y}}$, and $\mathring{y}_{k,i} = y_{k,i} - \bar{y}_i$ the demeaned return of asset $i$ in realization $k$. We
suppress the dependence on $t$ of the scalar quantities that follow. The sample covariance matrix $\bar{\bm{\Sigma}}_t$
has the entries
\begin{equation*}
    \bar{\sigma}_{ij} = \frac{1}{q} \sum_{k=1}^q \mathring{y}_{k,i}\, \mathring{y}_{k,j},
\end{equation*}
from which we obtain the sample correlations $r_{ij} = \bar{\sigma}_{ij} / \sqrt{\bar{\sigma}_{ii} \bar{\sigma}_{jj}}$
and their average $\bar{r} = \frac{2}{N(N-1)} \sum_{i<j} r_{ij}$. The constant-correlation target $\bm{T}_t$ has the
diagonal entries $\bar{\sigma}_{ii}$ and the off-diagonal entries $\bar{r} \sqrt{\bar{\sigma}_{ii} \bar{\sigma}_{jj}}$,
and the Ledoit--Wolf estimate is the convex combination
\begin{equation*}
    \hat{\bm{\Sigma}}_t = \hat{\delta}\, \bm{T}_t + (1 - \hat{\delta})\, \bar{\bm{\Sigma}}_t
\end{equation*}
with the shrinkage intensity $\hat{\delta} \in [0, 1]$.

Ledoit and Wolf choose the intensity that minimizes the expected squared Frobenius distance between
$\hat{\bm{\Sigma}}_t$ and $\bm{\Sigma}_t$, that is, the expected sum of the squared differences between their entries.
The optimal intensity is asymptotically proportional to $1/q$, and its estimate is
\begin{equation*}
    \hat{\delta} = \max\Bigl\{ 0, \min\Bigl\{ \frac{\hat{\pi} - \hat{\nu}}{q\, \hat{\omega}}, 1 \Bigr\} \Bigr\},
\end{equation*}
where
\begin{align*}
    \hat{\pi} &= \sum_{i,j} \hat{\pi}_{ij},
    \qquad \hat{\pi}_{ij} = \frac{1}{q} \sum_{k=1}^q \bigl( \mathring{y}_{k,i}\, \mathring{y}_{k,j} - \bar{\sigma}_{ij} \bigr)^2, \\
    \hat{\nu} &= \sum_{i} \hat{\pi}_{ii} + \bar{r} \sum_{i \neq j} \sqrt{\bar{\sigma}_{jj} / \bar{\sigma}_{ii}}\;
    \hat{\vartheta}_{ij}, \\
    \hat{\vartheta}_{ij} &= \frac{1}{q} \sum_{k=1}^q \bigl( \mathring{y}_{k,i}^2 - \bar{\sigma}_{ii} \bigr)
    \bigl( \mathring{y}_{k,i}\, \mathring{y}_{k,j} - \bar{\sigma}_{ij} \bigr), \\
    \hat{\omega} &= \sum_{i \neq j} \bigl( \bar{r} \sqrt{\bar{\sigma}_{ii} \bar{\sigma}_{jj}} - \bar{\sigma}_{ij} \bigr)^2.
\end{align*}
The quantity $\hat{\pi}$ estimates the sum of the asymptotic variances of the entries of $\sqrt{q}\,\bar{\bm{\Sigma}}_t$,
and $\hat{\nu}$ the sum of the asymptotic covariances between the entries of $\sqrt{q}\,\bm{T}_t$ and those of
$\sqrt{q}\,\bar{\bm{\Sigma}}_t$, which arise because both matrices are computed from the same realizations. The quantity
$\hat{\omega}$ is the squared Frobenius distance between $\bm{T}_t$ and $\bar{\bm{\Sigma}}_t$, which measures how far the
structure of the target is from the data. The intensity therefore increases with the estimation error of the sample
covariance matrix, decreases with the misspecification of the target, and tends to zero as $q$ grows, so that
$\hat{\bm{\Sigma}}_t$ approaches the sample covariance matrix when many observations are available. For the positive
average correlations typical of stock returns, the target is positive definite, and $\hat{\bm{\Sigma}}_t$ is then
positive definite whenever $\hat{\delta} > 0$, even when $N \ge q$.

The variance program replaces $\bm{\Sigma}_t$ by the Ledoit--Wolf estimate,
\begin{equation*} \label{eq:variance-program}
    \bm{x}_t^* = \argmin_{\bm{x}_t} \bm{x}_t^\top \hat{\bm{\Sigma}}_t \bm{x}_t, \quad \text{s.t. } \bm{x}_t^\top \bm{1}_N = 1,
\end{equation*}
and we write it in the standard form of a quadratic program in Appendix~\ref{app:mv_var_qp}.

One of the common criticisms of the variance is that it penalizes returns above the expected portfolio return as much as returns below it, although an
investor is typically concerned with underperformance rather than outperformance \cite{Markowitz1993}. This motivates
the semi-variance, which we consider next. 

\subsection{Semi-Variance}\label{sec:risk-semi-variance}

The semi-variance only considers the deviations of the portfolio return
below a threshold, so that returns above the threshold do not add to the risk of the portfolio. Markowitz et
al.~\cite{Markowitz1993} consider two thresholds, the expected portfolio return and a benchmark return, which may also be
a constant. Measured against a benchmark, the semi-variance captures the risk of underperforming the benchmark, while
measured against the expected return, it captures the dispersion of the portfolio return below its mean. We use the
expected portfolio return as the threshold. With $z^+ = \max\{z, 0\}$, the semi-variance of the portfolio return is
\begin{equation*}
    \text{SVar}_t\bigl( \bm{x}_t^\top \bm{R}_{t+1} \bigr)
    = \mathbb{E}_t\Bigl[ \bigl( (\bm{x}_t^\top \bm{\mu}_t - \bm{x}_t^\top \bm{R}_{t+1})^+ \bigr)^2 \Bigr].
\end{equation*}

As for the variance, the
conditional expectation has to be estimated. With the historical realizations $\bm{y}_1, \dots, \bm{y}_q$ and their
sample mean $\bar{\bm{y}}$ of \Cref{sec:risk-variance}, we replace $\bm{\mu}_t$ by $\bar{\bm{y}}$ and the conditional
expectation by the empirical average over the realizations, which gives the sample semi-variance
\begin{equation}\label{eq:sample-semi-variance}
    \frac{1}{q} \sum_{k=1}^q \bigl( (\bm{x}_t^\top \bar{\bm{y}} - \bm{x}_t^\top \bm{y}_k)^+ \bigr)^2
    = \bm{x}_t^\top \bm{\Sigma}^-(\bm{x}_t)\, \bm{x}_t,
\end{equation}
where
\begin{equation*}
    \bm{\Sigma}^-(\bm{x}_t) = \frac{1}{q} \sum_{k \,:\, \bm{x}_t^\top (\bm{y}_k - \bar{\bm{y}}) < 0}
    (\bm{y}_k - \bar{\bm{y}})(\bm{y}_k - \bar{\bm{y}})^\top
\end{equation*}
is the sample semi-covariance matrix of the realizations in which the portfolio return falls below its sample mean.

The sample semi-covariance matrix has the same form as the sample covariance matrix of \Cref{sec:risk-variance}, but it
only contains the realizations in which the portfolio underperforms its mean, and which realizations these are depends
on the weights. The semi-covariance matrix is therefore endogenous \cite{Estrada2008}, so unlike for the variance,
there is no fixed matrix that could be estimated before the optimization. For the same reason, the problem has no
closed-form solution, and the Ledoit--Wolf estimator of \Cref{sec:risk-variance} cannot be used. Estrada \cite{Estrada2008} proposes a heuristic semi-covariance
matrix that does not depend on the weights, but it only approximates the semi-variance of the portfolio. We therefore
minimize the sample semi-variance directly.

Each summand of the sample semi-variance \eqref{eq:sample-semi-variance} is only piecewise quadratic in the weights. Following Markowitz et al.~\cite{Markowitz1993}, we introduce for each realization $k$ an
auxiliary variable $u_k \ge 0$ subject to $u_k \ge \bm{x}_t^\top \bar{\bm{y}} - \bm{x}_t^\top \bm{y}_k$, and collect
them in $\bm{u} = [u_1, \dots, u_q]^\top$. Since the objective $\frac{1}{q} \sum_{k=1}^q u_k^2$ increases in every
$u_k \ge 0$, each $u_k$ equals $(\bm{x}_t^\top \bar{\bm{y}} - \bm{x}_t^\top \bm{y}_k)^+$ at the optimum, and the
semi-variance program is the quadratic program
\begin{equation*} 
\begin{aligned}
    \bm{x}_t^*, \bm{u}^* = \argmin_{\bm{x}_t,\, \bm{u}} \quad & \frac{1}{q} \sum_{k=1}^q u_k^2, \\
    \text{s.t.} \quad & \bm{x}_t^\top \bm{1}_N = 1, \\
    & \bm{x}_t^\top (\bm{y}_k - \bar{\bm{y}}) + u_k \ge 0, \quad u_k \ge 0, \quad k = 1, \dots, q.
\end{aligned}
\end{equation*}
Each realization adds one auxiliary variable and one linear constraint, so that the size of the program grows with the
number of realizations, whereas the variance program works with a covariance matrix whose size grows quadratically with
the number of assets. We write the program in the standard form of a quadratic program in
Appendix~\ref{app:mv_svar_qp}.

Like the variance, the semi-variance squares the deviations, so that a few large deviations below the mean can dominate
its estimate, and minimizing it requires a quadratic program. The mean absolute deviation, which we consider next,
measures the deviations linearly instead.

\subsection{Mean Absolute Deviation}\label{sec:risk-mad}

Konno and Yamazaki \cite{KonnoYamazaki1991} propose the mean absolute deviation as an alternative to the variance. It
measures the deviations of the portfolio return from its expected return by their absolute values instead of their
squares. Large deviations therefore have less
influence on the measure. The mean absolute deviation of the portfolio return is
\begin{equation*}
    \text{MAD}_t\bigl( \bm{x}_t^\top \bm{R}_{t+1} \bigr)
    = \mathbb{E}_t\bigl[ \lvert \bm{x}_t^\top \bm{\mu}_t - \bm{x}_t^\top \bm{R}_{t+1} \rvert \bigr].
\end{equation*}
The main motivation of Konno and Yamazaki is computational. Since the absolute value is piecewise linear, minimizing the
sample mean absolute deviation can be formulated as a linear program rather than a quadratic program, which makes it more efficient for larger problems, 
and setting up the problem requires no estimate of the covariance matrix.

As for the semi-variance, we replace $\bm{\mu}_t$ by the sample mean $\bar{\bm{y}}$ and the conditional expectation by
the empirical average over the realizations $\bm{y}_1, \dots, \bm{y}_q$ of \Cref{sec:risk-variance}, which gives the
sample mean absolute deviation $\frac{1}{q} \sum_{k=1}^q \lvert \bm{x}_t^\top \bar{\bm{y}} - \bm{x}_t^\top \bm{y}_k
\rvert$. Since the deviations $\bm{x}_t^\top \bar{\bm{y}} - \bm{x}_t^\top \bm{y}_k$ sum to zero over the sample, the sum
of their positive parts equals the sum of their negative parts, and the absolute value of each deviation is the sum of
its two parts. The sample mean absolute deviation is therefore twice the average shortfall of the portfolio return below
its sample mean \cite{Speranza1993, OgryczakRuszczynski1999}
\begin{equation*}
    \frac{1}{q} \sum_{k=1}^q \lvert \bm{x}_t^\top \bar{\bm{y}} - \bm{x}_t^\top \bm{y}_k \rvert
    = \frac{2}{q} \sum_{k=1}^q (\bm{x}_t^\top \bar{\bm{y}} - \bm{x}_t^\top \bm{y}_k)^+.
\end{equation*}

To write the problem as a linear program, we introduce, as for the semi-variance, an auxiliary variable $u_k \ge 0$ for
each realization $k$ subject to $u_k \ge \bm{x}_t^\top \bar{\bm{y}} - \bm{x}_t^\top \bm{y}_k$.
Since the objective $\frac{2}{q} \sum_{k=1}^q u_k$ increases in every $u_k$, each $u_k$ equals
$(\bm{x}_t^\top \bar{\bm{y}} - \bm{x}_t^\top \bm{y}_k)^+$ at the optimum, and the mean absolute deviation program is the
linear program
\begin{equation*}
\begin{aligned}
    \bm{x}_t^*, \bm{u}^* = \argmin_{\bm{x}_t,\, \bm{u}} \quad & \frac{2}{q} \sum_{k=1}^q u_k, \\
    \text{s.t.} \quad & \bm{x}_t^\top \bm{1}_N = 1, \\
    & \bm{x}_t^\top (\bm{y}_k - \bar{\bm{y}}) + u_k \ge 0, \quad u_k \ge 0, \quad k = 1, \dots, q.
\end{aligned}
\end{equation*}
It has the same constraints as the semi-variance program, with the auxiliary variables in the objective in linear form. We write it in the standard form of a linear program in Appendix~\ref{app:mad_representation}.

The linear program is computationally efficient, but its solution is sensitive to estimation error. Kondor et
al.~\cite{kondor2007noise} show that portfolios optimized under the mean absolute deviation are even more sensitive to
estimation error than those optimized under the variance, and that their weights fluctuate more strongly from sample to
sample.

The three risk measures considered so far measure the dispersion of the portfolio return around its mean. Expected
Shortfall, which we consider next, instead measures the size of the largest losses.

\subsection{Expected Shortfall}\label{sec:risk-es}

Expected Shortfall measures the risk of a portfolio by the average of its largest losses. Let
$L(\bm{x}_t) = -\bm{x}_t^\top \bm{R}_{t+1}$ denote the loss of the portfolio at time $t+1$ and $\mathbb{P}_t(\cdot)$
the probability conditional on the information available at time $t$. The Value-at-Risk (VaR) at confidence level
$\beta \in (0, 1)$ is the $\beta$-quantile of the loss distribution,
\begin{equation*}
    \text{VaR}_\beta(\bm{x}_t) = \min \{ \alpha \mid \mathbb{P}_t(L(\bm{x}_t) \le \alpha) \ge \beta \}.
\end{equation*}
VaR does not take into account how large the losses beyond it are. Expected Shortfall is the expected loss given that
the loss exceeds VaR,
\begin{equation*}
    \text{ES}_\beta(\bm{x}_t) = \mathbb{E}_t\bigl[ L(\bm{x}_t) \mid L(\bm{x}_t) \ge \text{VaR}_\beta(\bm{x}_t) \bigr],
\end{equation*}
so that it accounts for both the probability and the size of the losses in the tail. For loss distributions with atoms,
such as the empirical distribution used below, Expected Shortfall is defined as the average of the quantiles of the loss
distribution above $\beta$, which is equivalent to the conditional expectation for continuous distributions
\cite{acerbi2002coherenceexpectedshortfall}. In its capital requirements for market risk, the Basel Committee on Banking
Supervision has replaced VaR at the $99\%$ level by Expected Shortfall at the $97.5\%$ level \cite{BCBS2019}, indicating
that it is a more appropriate measure of risk than VaR.

Since $\text{VaR}_\beta(\bm{x}_t)$ itself depends on the portfolio weights, the conditional expectation in the definition is not
in a suitable form for optimization. Rockafellar and Uryasev \cite{RockafellarUryasev2000, RockafellarUryasev2002} show
that Expected Shortfall is the minimum over $\alpha \in \R$ of the auxiliary function
\begin{equation*}
    F_\beta(\bm{x}_t, \alpha) = \alpha + \frac{1}{1-\beta}\, \mathbb{E}_t\bigl[ (L(\bm{x}_t) - \alpha)^+ \bigr],
\end{equation*}
and that the minimizing $\alpha$ is $\text{VaR}_\beta(\bm{x}_t)$. Minimizing Expected Shortfall over the portfolio weights $\bm{x}_t$ is
therefore equivalent to minimizing $F_\beta$ jointly over $\bm{x}_t$ and $\alpha$,
\begin{equation*}
    \min_{\bm{x}_t} \text{ES}_\beta(\bm{x}_t) = \min_{\bm{x}_t, \alpha} F_\beta(\bm{x}_t, \alpha),
\end{equation*}
where $F_\beta$ is jointly convex in $(\bm{x}_t, \alpha)$. The derivation of this representation is given in Appendix~\ref{app:es_auxiliary}.

With the historical realizations $\bm{y}_1, \dots, \bm{y}_q$ of \Cref{sec:risk-variance}, we replace the conditional
expectation in $F_\beta$ by the empirical average, which gives the objective
$\alpha + \frac{1}{q(1-\beta)} \sum_{k=1}^q (-\bm{x}_t^\top \bm{y}_k - \alpha)^+$. As for the semi-variance and the
mean absolute deviation, we introduce for each realization $k$ an auxiliary variable $u_k \ge 0$ subject to
$u_k \ge -\bm{x}_t^\top \bm{y}_k - \alpha$. Since the objective increases in every $u_k$, each $u_k$ equals
$(-\bm{x}_t^\top \bm{y}_k - \alpha)^+$ at the optimum, and the Expected Shortfall program is the linear program
\begin{equation*}
\begin{aligned}
    \bm{x}_t^*, \alpha^*, \bm{u}^* = \argmin_{\bm{x}_t,\, \alpha,\, \bm{u}} \quad
    & \alpha + \frac{1}{q(1-\beta)} \sum_{k=1}^q u_k, \\
    \text{s.t.} \quad & \bm{x}_t^\top \bm{1}_N = 1, \\
    & \bm{x}_t^\top \bm{y}_k + \alpha + u_k \ge 0, \quad u_k \ge 0, \quad k = 1, \dots, q.
\end{aligned}
\end{equation*}
Since Expected Shortfall measures losses rather than deviations from the mean, the constraints contain the realizations
$\bm{y}_k$ themselves instead of their deviations from the sample mean. We write the program in the standard form of a
linear program in Appendix~\ref{app:lp_es}.

At the optimum, only the realizations whose loss exceeds $\alpha^*$, which are about the $(1-\beta) q$ realizations with
the largest losses, contribute to the objective. The estimate of Expected Shortfall is therefore based on a fraction of the
sample, and Kondor et al.~\cite{kondor2007noise} show that portfolios optimized under Expected Shortfall are more
sensitive to estimation error than those optimized under the variance and the mean absolute deviation. Goldberg et
al.~\cite{goldberg2013minimizingshortfall} note that estimation error matters more for Expected Shortfall than for the
variance, since most realizations only serve to determine which losses are the largest. The confidence level thus
trades off the focus on the tail against the number of realizations that determine the estimate. We use $\beta = 0.8$,
at which about $500$ of the $q \approx 2{,}500$ daily realizations contribute to the objective. Testing confidence levels from $60\%$ to $99\%$, Goldberg et al.~\cite{goldberg2013minimizingshortfall} also find
that minimum Expected Shortfall portfolios outperform minimum-variance portfolios most at relatively low confidence
levels, at which Expected Shortfall captures the asymmetry of the return distribution rather than its extreme losses.

Furthermore, Ciliberti et
al.~\cite{ciliberti2007feasibility} show that, without bounds on the portfolio weights, the program is unbounded in some samples. 
For large portfolios, it
has no finite solution once $N/q$ exceeds a critical value, which is about $0.5$ for high confidence levels, and the weights fluctuate from sample to sample as $N/q$ approaches this value. In the large
universe of the empirical analysis, which contains up to $1{,}400$ stocks, $N/q$ can come close to it. The leverage
constraint that we add later in this chapter bounds the positions and thereby ensures that the program has a solution.

Expected Shortfall weighs all losses beyond VaR equally, so that the largest losses in the tail count no more
than the losses just beyond the quantile. Second-order Expected Shortfall, which we consider next, puts more weight on
the largest losses.

\subsection{Second-Order Expected Shortfall}\label{sec:risk-soes}

Second-order Expected Shortfall belongs to the family of higher-moment coherent risk measures of Krokhmal
\cite{Krokhmal2007}, which measure the risk in the tail of the loss distribution through its higher moments and contain
Expected Shortfall as the first-order member. The second-order member replaces the expected excess loss in the auxiliary
function $F_\beta$ of \Cref{sec:risk-es} by its $L^2$ norm $\|Z\|_{L^2} = (\mathbb{E}_t[Z^2])^{1/2}$, which gives the
auxiliary function
\begin{equation*}
    G_\beta(\bm{x}_t, \alpha) = \alpha + \frac{1}{1-\beta}\, \bigl\| (L(\bm{x}_t) - \alpha)^+ \bigr\|_{L^2},
\end{equation*}
and second-order Expected Shortfall at confidence level $\beta$ is defined as its minimum over $\alpha$,
\begin{equation*}
    \text{SOES}_\beta(\bm{x}_t) = \min_{\alpha \in \R} G_\beta(\bm{x}_t, \alpha).
\end{equation*}
Since the $L^2$ norm squares the excess losses, large losses beyond $\alpha$ receive more weight than in Expected
Shortfall. Krokhmal notes that second-order Expected Shortfall measures the risk in the tail through the second moment
of the losses, which relates it to the mean-variance framework. In a case study on fifty S\&P~500 stocks with
heavy-tailed returns, he finds that portfolios minimizing the higher-moment measures at the $90\%$ level outperform
those minimizing Expected Shortfall or the variance out of sample when the required expected return is low, while the
strategies perform similarly when it is high, and he stresses that the results are specific to the data.

By definition, minimizing second-order Expected Shortfall is
\begin{equation*}
    \min_{\bm{x}_t} \text{SOES}_\beta(\bm{x}_t) = \min_{\bm{x}_t, \alpha} G_\beta(\bm{x}_t, \alpha),
\end{equation*}
where $G_\beta$ is jointly convex in $(\bm{x}_t, \alpha)$, which we show in Appendix~\ref{app:soes_auxiliary}.

With the historical realizations $\bm{y}_1, \dots, \bm{y}_q$ of \Cref{sec:risk-variance}, we replace the $L^2$ norm in
$G_\beta$ by its empirical counterpart
$\bigl( \frac{1}{q} \sum_{k=1}^q \bigl( (-\bm{x}_t^\top \bm{y}_k - \alpha)^+ \bigr)^2 \bigr)^{1/2}$. As for Expected
Shortfall, we introduce for each realization $k$ an auxiliary variable $u_k \ge 0$ subject to
$u_k \ge -\bm{x}_t^\top \bm{y}_k - \alpha$, so that the empirical norm becomes $\frac{1}{\sqrt{q}} \|\bm{u}\|_2$ and each
$u_k$ equals $(-\bm{x}_t^\top \bm{y}_k - \alpha)^+$ at the optimum. Since the objective is still not linear, we introduce
one more auxiliary variable $s$ subject to $\|\bm{u}\|_2 \le s$, which equals $\|\bm{u}\|_2$ at the optimum, and the
second-order Expected Shortfall program is
\begin{equation*}
\begin{aligned}
    \bm{x}_t^*, \alpha^*, \bm{u}^*, s^* = \argmin_{\bm{x}_t,\, \alpha,\, \bm{u},\, s} \quad
    & \alpha + \frac{1}{(1-\beta)\sqrt{q}}\, s, \\
    \text{s.t.} \quad & \bm{x}_t^\top \bm{1}_N = 1, \\
    & \bm{x}_t^\top \bm{y}_k + \alpha + u_k \ge 0, \quad u_k \ge 0, \quad k = 1, \dots, q, \\
    & \|\bm{u}\|_2 \le s.
\end{aligned}
\end{equation*}
The constraint $\|\bm{u}\|_2 \le s$ requires the vector $[\bm{u}^\top, s]^\top$ to lie in the second-order cone, the set
of vectors whose last entry is at least the Euclidean norm of the remaining entries \cite[Section~2.2.3]{boyd2004convex},
so that the program is a second-order cone program \cite[Section~4.4.2]{boyd2004convex}. We write it in standard form in
Appendix~\ref{app:socp_soes}.

The cone constraint makes the program harder to solve than the linear and quadratic programs of the other risk
measures, and the second-order Expected Shortfall programs are computationally the most expensive programs in our
backtests (Appendix~\ref{sec:computational-complexity}). We use the same confidence level $\beta = 0.8$ as for Expected
Shortfall, so that the two tail measures can be compared, and as for Expected Shortfall, the leverage constraint that we
add later in this chapter bounds the positions and thereby ensures that the program has a solution.

Expected Shortfall and second-order Expected Shortfall measure the size of the largest losses, whereas the variance,
the semi-variance, and the mean absolute deviation measure deviations from the mean. We compare these two classes of
risk measures next.

\subsection{Coherent Risk Measures and Deviation Measures}\label{sec:coherent-and-deviation-measures}

The five risk measures of this section fall into two classes, which differ in how they treat the location of the return
distribution. To describe them, we write $X$ for the portfolio return $\bm{x}_t^\top \bm{R}_{t+1}$ and $C$ for a
constant.

Artzner et al.~\cite{CoherentMeasuresOfRisk} call a risk measure $\mathcal{R}$ coherent if it satisfies four
properties. Adding a riskless amount to the return reduces the risk by that amount,
$\mathcal{R}(X + C) = \mathcal{R}(X) - C$, scaling the position by a positive factor scales the risk by the same
factor, the risk of a combined position is at most the sum of the risks of its parts, so that diversification does not
increase risk, and a return that is higher in every state of the world carries no more risk. VaR violates the third
property, subadditivity, in general, whereas Expected Shortfall \cite{acerbi2002coherenceexpectedshortfall} and
second-order Expected Shortfall \cite{Krokhmal2007} are coherent. Both are moreover strictly expectation bounded, that
is, their risk exceeds the expected loss of every nonconstant return.

Rockafellar et al.~\cite{RockafellarUryasevZabarankin2006} define deviation measures $\mathcal{D}$, which quantify
the uncertainty of a return rather than its losses. They share positive homogeneity and subadditivity with coherent risk
measures, are non-negative and zero only for constant returns, and replace the first property by translation
invariance, $\mathcal{D}(X + C) = \mathcal{D}(X)$. The standard deviation, the semi-deviation, which is the square root
of the semi-variance, and the mean absolute deviation are deviation measures. The variance and the semi-variance are the
squares of the standard deviation and the semi-deviation, so they are not positively homogeneous and therefore not
deviation measures themselves. They are, however, translation invariant, and since squaring is increasing on
non-negative numbers, they have the same minimizers as the standard deviation and the semi-deviation.

The two translation properties capture the difference between the classes. A deviation measure depends only on the
dispersion of the return around its mean and is unaffected by a shift of the whole distribution, whereas a coherent
risk measure decreases when the distribution shifts upward. We state the formal definitions of both classes and prove
that Expected Shortfall and second-order Expected Shortfall are coherent and strictly expectation bounded in
Appendix~\ref{app:demeaning}.

So far, we have treated the asset returns as given and denominated in the investor's domestic currency. An
international investor, however, holds assets denominated in several currencies, whose returns in the domestic currency
also depend on the exchange rates and on the currency positions that the investor takes to hedge them. All five risk
measures are evaluated on the portfolio return in the domestic currency, so before we can formulate the multi-currency
programs, we have to express this return in terms of the asset weights and the currency positions. The next section
derives this representation exactly, without assuming a model for the joint distribution of asset and currency returns.

\section{Model-Free Multi-Currency Portfolio Returns}\label{sec:returns}
\subsection{Unhedged Portfolio Returns}

We consider an international investor who allocates their wealth across $N$ assets denominated in up to $M+1$ currencies. We fix a discrete time index $t\in\{0,1,2,\dots\}$ and one investment period from $t$ to $t{+}1$. Let the investor's \emph{domestic currency} be indexed by $0$ and let the remaining $M$ foreign currencies be indexed by $\mathcal{C} \coloneqq \{1,2,\dots,M\}$.

Let $P_{i,t}>0$ be the price of asset $i$ at time $t$ in its local currency $c_i \in \mathcal{C}\cup \{0\}$. Define the corresponding \emph{simple local-currency asset return} over $(t,t{+}1]$ as
\begin{equation*}
R_{i,t+1} \coloneqq \frac{P_{i,t+1}-P_{i,t}}{P_{i,t}}.
\end{equation*}
Thus $P_{i,t+1}=P_{i,t}(1+R_{i,t+1})$.

For each foreign currency $c\in\mathcal{C}$, let $S_{c,t} > 0$ denote the exchange rate at time $t$, quoted as units of domestic currency per one unit of foreign currency $c$. The associated \emph{simple currency return} over $(t,t{+}1]$ is
\begin{equation*}
e_{c,t+1} \coloneqq \frac{S_{c,t+1}-S_{c,t}}{S_{c,t}}.
\end{equation*}
Equivalently, $S_{c,t+1} = S_{c,t}(1+e_{c,t+1})$. For the domestic currency, we set $S_{0,t}\equiv 1$ and $e_{0,t+1}\equiv 0$.

The domestic-currency value of holding one unit of asset $i$ at time $t$ is $P_{i,t}\,S_{c_i,t}$. Therefore, the \emph{unhedged domestic-currency simple return} of asset $i$ is
\begin{align}
R^{u}_{i,t+1}
&=
\frac{P_{i,t+1}S_{c_i,t+1} - P_{i,t}S_{c_i,t}}{P_{i,t}S_{c_i,t}}
\nonumber\\
&=
\frac{P_{i,t}(1+R_{i,t+1})\,S_{c_i,t}(1+e_{c_i,t+1}) - P_{i,t}S_{c_i,t}}{P_{i,t}S_{c_i,t}}
\nonumber\\
&=
R_{i,t+1} + e_{c_i,t+1} + R_{i,t+1}\,e_{c_i,t+1}.
\label{eq:unhedged-asset-return}
\end{align}

Let $\bm{x}_t\in\R^N$ be the vector of asset portfolio weights chosen at time $t$ and held until $t{+}1$, where $x_{i,t}$ is the fraction of total portfolio value invested in asset $i$ at time $t$. We impose the standard budget constraint $\bm{1}_N^\top \bm{x}_t = 1$.

Assuming no rebalancing between times $t$ and $t{+}1$, the domestic-currency portfolio's unhedged return is given by
\begin{equation*}
R^{u}_{\mathcal{P},t+1} = \sum_{i=1}^N x_{i,t}\,R^{u}_{i,t+1}.
\end{equation*}

\subsection{Hedged Portfolio Returns}

To manage currency risk, the investor can take positions in one-period currency forward contracts. For each foreign currency $c\in\mathcal{C}$, let $F_{c,t}$ be the forward rate quoted at time $t$ for settlement at $t{+}1$, in units of domestic currency per unit of foreign currency $c$. We define the \emph{forward premium} by
\begin{equation*}
f_{c,t} \coloneqq \frac{F_{c,t}-S_{c,t}}{S_{c,t}}.
\end{equation*}

The investor hedges by selling foreign currency forward. Let $\upphi_{c,t}$ denote the relative notional of the \emph{short} forward position in currency $c$ at time $t$, expressed as a fraction of total portfolio value, with the convention that $\upphi_{c,t}>0$ corresponds to hedging. Per unit of notional, the short position pays $f_{c,t}-e_{c,t+1}$ in domestic currency over the period. Adding these forward payoffs to the unhedged asset returns, the \emph{hedged portfolio return} is
\begin{equation}
R^{h}_{\mathcal{P},t+1}
=
\underbrace{\sum_{i=1}^N x_{i,t}\,R^{u}_{i,t+1}}_{R^{u}_{\mathcal{P},t+1}}
+
\sum_{c=1}^M \upphi_{c,t}\,(f_{c,t}-e_{c,t+1}).
\label{eq:hedged-portfolio-return}
\end{equation}

\subsection{Matrix Representation and Decomposition}\label{sec:matrix-representation}

To express the returns in a compact form suitable for optimization, we define the currency incidence matrix $\bm{C}\in\{0,1\}^{N\times M}$ by
\begin{equation*}
C_{i,c} \coloneqq
\begin{cases}
1, & \text{if asset $i$ is denominated in foreign currency $c$},\\
0, & \text{otherwise}.
\end{cases}
\end{equation*}
Each row of $\bm{C}$ has at most one nonzero entry, and the rows corresponding to domestic assets are identically zero. This allows us to define the portfolio's \emph{implicit foreign-currency exposure} as
\begin{equation*}
\bm{w}_t \coloneqq \bm{C}^\top \bm{x}_t \in \R^M,
\qquad
w_{c,t} = \sum_{i:\,c_i=c} x_{i,t},
\end{equation*}
such that $w_{c,t}$ is the total portfolio weight invested in assets denominated in currency $c$. The implicit domestic exposure follows residually as $w_{0,t} = 1 - \bm{1}_M^\top \bm{w}_t$, since the asset weights sum to one.

We express returns in vector form $\bm{R}_{t+1} \in \R^N$, $\bm{e}_{t+1} \in \R^M$, and the forward premiums as $\bm{f}_t \in \R^M$. The product $\bm{C}\bm{e}_{t+1}\in\R^N$ aligns currency returns to assets, with $i$th component $e_{c_i,t+1}$ for foreign assets and $0$ for domestic ones. Using \eqref{eq:unhedged-asset-return}, the vector of \emph{unhedged domestic returns} across assets can be written as
\begin{equation*}
\bm{R}^{u}_{t+1}
=
\bm{R}_{t+1}
+
\bm{C}\bm{e}_{t+1}
+
(\bm{R}_{t+1}\had\bm{C}\bm{e}_{t+1}),
\end{equation*}
where $\had$ denotes the Hadamard product. 

Expressing the forward positions in vector form as $\bm{\upphi}_t\in\R^M$, the hedged portfolio return \eqref{eq:hedged-portfolio-return} can be written as
\begin{equation}
R^{h}_{\mathcal{P},t+1} = \bm{x}_t^\top \bm{R}^{u}_{t+1} + \bm{\upphi}_t^\top(\bm{f}_t-\bm{e}_{t+1}).
\label{eq:hedged-vector}
\end{equation}

Setting $\bm{\upphi}_t=\bm{w}_t=\bm{C}^\top\bm{x}_t$ neutralizes the implicit currency exposure. Substituting this together with the unhedged return vector into \eqref{eq:hedged-vector}, the \emph{fully hedged portfolio return} is
\begin{align*}
R^{fh}_{\mathcal{P},t+1}
&= \bm{x}_t^\top\bm{R}^{u}_{t+1} + \big(\bm{C}^\top\bm{x}_t\big)^\top(\bm{f}_t-\bm{e}_{t+1})\\[2pt]
&= \bm{x}_t^\top\Big[\bm{R}_{t+1} + \bm{C}\bm{e}_{t+1} + \big(\bm{R}_{t+1}\had\bm{C}\bm{e}_{t+1}\big)\Big] + \bm{x}_t^\top\bm{C}(\bm{f}_t-\bm{e}_{t+1})\\[2pt]
&= \bm{x}_t^\top\Big[\bm{R}_{t+1} + \bm{C}\bm{e}_{t+1} + \big(\bm{R}_{t+1}\had\bm{C}\bm{e}_{t+1}\big) + \bm{C}\bm{f}_t - \bm{C}\bm{e}_{t+1}\Big]\\[2pt]
&= \bm{x}_t^\top\Big[\bm{R}_{t+1} + \bm{C}\bm{f}_t + \big(\bm{R}_{t+1}\had\bm{C}\bm{e}_{t+1}\big)\Big]\\[2pt]
&= \bm{x}_t^\top\bm{R}^{fh}_{t+1}.
\end{align*}
The first-order currency terms $\pm\,\bm{x}_t^\top\bm{C}\bm{e}_{t+1}$ cancel, leaving the \emph{fully hedged asset return vector}
\begin{equation*}
\bm{R}^{fh}_{t+1}
=
\bm{R}_{t+1}
+
\bm{C}\bm{f}_t
+
(\bm{R}_{t+1}\had\bm{C}\bm{e}_{t+1}).
\end{equation*}

More generally, define the \emph{net currency exposure} $\bm{\psi}_t \coloneqq \bm{w}_t-\bm{\upphi}_t$. Imposing $\sum_{c=0}^M \upphi_{c,t}=1$ gives $\sum_{c=0}^M \psi_{c,t}=0$, so the currency overlay is a zero-investment portfolio. Substituting $\bm{\upphi}_t=\bm{w}_t-\bm{\psi}_t$ into \eqref{eq:hedged-vector} expresses the hedged return in terms of net exposures,
\begin{align}
R^{h}_{\mathcal{P},t+1}
&= \bm{x}_t^\top \bm{R}^{u}_{t+1} + (\bm{w}_t-\bm{\psi}_t)^\top(\bm{f}_t-\bm{e}_{t+1})\nonumber \\[2pt]
&= \bm{x}_t^\top \bm{R}^{u}_{t+1} + \bm{w}_t^\top(\bm{f}_t-\bm{e}_{t+1}) + \bm{\psi}_t^\top(\bm{e}_{t+1}-\bm{f}_t) \nonumber\\[2pt]
&= \bm{x}_t^\top \bm{R}^{fh}_{t+1} + \bm{\psi}_t^\top(\bm{e}_{t+1}-\bm{f}_t),
\label{eq:hedged-net}
\end{align}

which we adopt throughout the following portfolio optimization formulations, as it makes the asset weights $\bm{x}_t$ and the net currency exposures $\bm{\psi}_t$ the explicit variables of the optimization.

\section{Regularized Multi-Currency Portfolio Optimization}\label{sec:regularized-multi-currency}

\subsection{Regularization}\label{sec:regularization}

\subsection{Sequential Optimization}\label{sec:sequential-optimization}

\subsection{Joint Optimization}\label{sec:joint-optimization}

\section{Mean-Risk Portfolio Optimization} \label{sec:mean-risk-portfolio-optimization}
